% TOP_PROBLEM: {"id":30007516,"problem_number":"LOCAL-30007516","title":"Self-fulfilling fluctuations","category_id":21,"category":{"id":21,"name":"economics","display_name":"Economics","description":"Open economic theory statements, published research agendas, and proposed scoped specifications; question provenance and historical priority are qualified per record.","slug":"economics","order_index":21,"created_at":"2026-10-08T00:00:00Z"},"status":"research_question","external_url":null,"legacy_ids":[],"published":false,"statement":"Identify the probability that suppressing independently identified belief innovations would prevent a recession while keeping fundamental shocks and policy rules fixed.","economics":{"original_id":5,"field":"Business cycles and macro dynamics","kind":"identification","formalization_basis":"adapted_research_question","source_status":"research_agenda","priority_status":"earliest_found","priority_note":"This is the earliest explicit relevant statement verified in this audit, not a claim of original priority; older literature and earlier versions may contain antecedents.","earliest_verified_reference":{"authors":["Benjamin Moll"],"title":"Research Agenda: Benjamin Moll","year":2020,"url":"https://economicdynamics.org/research-agenda-moll2020/","doi":null,"locator":"Section 4.3, “Heterogeneity, Bubbles and Crashes”","evidence_summary":"The author says the root causes of large crises are not soundly understood and asks for quantitatively disciplined theories using balance-sheet and portfolio evidence. Belief-driven crisis attribution below is a restricted adaptation."},"current_reference":{"authors":["Benjamin Moll"],"title":"Research Agenda: Benjamin Moll","year":2020,"url":"https://economicdynamics.org/research-agenda-moll2020/","doi":null,"locator":"Section 4.3, “Heterogeneity, Bubbles and Crashes”","evidence_summary":"The author says the root causes of large crises are not soundly understood and asks for quantitatively disciplined theories using balance-sheet and portfolio evidence. Belief-driven crisis attribution below is a restricted adaptation."},"formalization_note":"The survey asks for disciplined theories of large crises. The intervention-based attribution estimand is a restricted adaptation, not a published canonical sunspot conjecture."},"statusQualification":"Published question or research agenda verified at the cited locator. The mathematical specification is adapted for this catalogue and includes additional assumptions; source publication does not establish first-ever priority or certify this exact model remains unsolved. The survey asks for disciplined theories of large crises. The intervention-based attribution estimand is a restricted adaptation, not a published canonical sunspot conjecture.","statusReviewedAt":"2026-10-08"}
% FIRST_POSTED: 2026-10-08
% ENTRY_KIND: statement_only
% !TeX program = lualatex
\documentclass[11pt]{article}
\usepackage{fontspec}
\usepackage[a4paper,margin=25mm]{geometry}
\usepackage{amsmath,amssymb,mathtools}
\usepackage{xurl}
\usepackage[hidelinks]{hyperref}
\newcommand{\sref}[2]{\hyperlink{source:#1:#2}{[S#2]}}
\setlength{\emergencystretch}{3em}
\begin{document}
% problemId:30007516
\section*{Self-fulfilling fluctuations}\label{30007516}
\subsection{Definitions and mathematical statement}
Fix a nonlinear incomplete-markets model with observed fundamental state \(f_t\), endogenous distribution \(\mu_t\), and a declared equilibrium-selection variable \(b_t\) representing beliefs or coordination. A belief innovation \(\eta_t\) is required to be orthogonal to innovations in the specified fundamental-information set, not merely to lagged output. Let equilibrium log output be \(y_t=g(f_t,\mu_t,b_t)\). Define a recession as \(y_{t+H}-y_t<-q\), for preselected horizon \(H\) and threshold \(q>0\).
For each observed recession episode, construct the counterfactual path \(y^{0}_t\) obtained by suppressing \(\eta_t\) while retaining the same fundamental innovations, initial distribution, and policy rule. The target is
\[
\Theta=\Pr\{y_{t+H}-y_t<-q,\;
y^{0}_{t+H}-y^{0}_t\ge -q\},
\]
the probability that a recession would be prevented under this intervention on beliefs. Determine an identified set for \(\Theta\) from forecast surveys, balance sheets, asset prices, and instruments with explicit exclusion restrictions. Because equilibrium selection is latent and beliefs can react to fundamentals, provide observationally equivalent explanations and falsifiable restrictions. Establish robustness to plausible selection rules and financial-friction specifications. Existence of sunspot equilibria in a particular model does not establish their quantitative importance in actual recessions.

\par\textbf{Formulation and scope.} The survey asks for disciplined theories of large crises. The intervention-based attribution estimand is a restricted adaptation, not a published canonical sunspot conjecture.

\par\textbf{Open-question provenance.} An explicit question or research agenda was verified at the source locator. This does not by itself certify that every later version remains unresolved \sref{30007516}{1}.

\par\textbf{Historical priority.} This is the earliest explicit relevant statement verified in this audit, not a claim of original priority; older literature and earlier versions may contain antecedents.

\subsection{Short English statement}
Identify the probability that suppressing independently identified belief innovations would prevent a recession while keeping fundamental shocks and policy rules fixed.

\subsection{Sources}
\begin{itemize}\raggedright
\item[\textnormal{[S1]}]\hypertarget{source:30007516:1}{} Benjamin Moll. 2020. Research Agenda: Benjamin Moll. Section 4.3, “Heterogeneity, Bubbles and Crashes”.
\url{https://economicdynamics.org/research-agenda-moll2020/}.
{\small\textit{Earliest verified question/agenda reference; historical priority qualified below. The author says the root causes of large crises are not soundly understood and asks for quantitatively disciplined theories using balance-sheet and portfolio evidence. Belief-driven crisis attribution below is a restricted adaptation.}}
\end{itemize}
\end{document}
